\begin{definition}[4.7.4]
\label{IV.2.4.7.4-fr}
Sean $k$ un cuerpo, $X$ un $k$-preesquema íntegro y $K$ el cuerpo de funciones
racionales de $X$. Se llama \emph{multiplicidad radicial} de $K$
(o de $X$) sobre $k$ al supremo de las longitudes de los anillos artinianos
$K\otimes_k k_1$, cuando $k_1$ recorre el conjunto de las extensiones
radiciales finitas de $k$. Se llama \emph{multiplicidad separable} de $K$ (o de
$X$) sobre $k$ al número geométrico de las componentes irreducibles de $X$ si
este número es finito, y $+\infty$ en caso contrario. Finalmente, se llama
\emph{multiplicidad total} de $K$ (o de $X$) sobre $k$ al producto de la
multiplicidad radicial y la multiplicidad separable de $K$ sobre $k$.
\end{definition}

Se observará que, si $p$ es el exponente característico de $k$ y si la
multiplicidad radicial de $K$ sobre $k$ es finita, es una potencia $q=p^f$ de
$p$, como se deduce de (4.7.3, (ii)); cuando $p\ne1$, $f$ está bien determinado
y se llama \emph{exponente de inseparabilidad} de $K$ sobre $k$; cuando $p=1$,
la multiplicidad radicial es siempre igual a $1$.

Decir que la multiplicidad radicial (resp. separable, resp. total) de $X$ sobre
$k$ es $1$ significa que $X$ es geométricamente reducido (resp.
geométricamente irreducible, resp. geométricamente íntegro) sobre $k$.

Cuando $X$ es localmente de tipo finito, la multiplicidad radicial de $X$ sobre
$k$ es también la longitud común de los anillos locales $A_i$ en los ideales
primos minimales del anillo total de fracciones $A$ de $K\otimes_k k'$, para
toda extensión perfecta $k'$ de $k$; la multiplicidad separable de $X$ sobre
$k$ es el número de ideales primos minimales de $A=K\otimes_k\Omega$ para toda
extensión algebraicamente cerrada $\Omega$ de $k$, y, finalmente, la
multiplicidad total es la suma de las longitudes de los anillos locales $A_i$
de $A=K\otimes_k\Omega$ en sus ideales primos minimales~; estos tres números
son \emph{finitos} en este caso.

\begin{definition}[4.7.5]
\label{IV.2.4.7.5-fr}
Sean $k$ un cuerpo, $X$ un $k$-preesquema localmente de tipo finito sobre $k$,
$\mathscr F$ un $\mathcal O_X$-módulo coherente y $x$ un punto de $X$ tal que
$\mathscr F_x$ sea un $\mathcal O_x$-módulo de longitud finita. Se llama
\emph{longitud geométrica} de $\mathscr F$ en $x$ (relativa a $k$), o
\emph{multiplicidad radicial} de $x$ para $\mathscr F$, al producto
$\lambda_{\boldsymbol{x}}(\mathscr F)
=(\operatorname{long}_{\mathcal O_{\boldsymbol{x}}}(\mathscr F_{\boldsymbol{x}}))\mu_i(\boldsymbol{k}(x)\mid k)$,
donde $\mu_i(\boldsymbol{k}(x)\mid k)$ designa la multiplicidad radicial de $\boldsymbol{k}(x)$ sobre
$k$.
\end{definition}

Cuando $X$ es \emph{íntegro} y $x$ es el punto genérico de $X$, se tiene
$\operatorname{long}(\mathcal O_x)=1$, por lo que la multiplicidad radicial de
$x$ para $\mathcal O_X$ no es otra que la multiplicidad radicial de $X$ sobre
$k$, definida en (4.7.4).

\oldpage[IV-2]{78}

\begin{proposition}[4.7.6]
\label{IV.2.4.7.6-fr}
Sea
$0\boldsymbol{\longrightarrow}\mathscr F'\longrightarrow\mathscr F
\longrightarrow\mathscr F''\longrightarrow0$
una sucesión exacta de $\mathcal O_X$-módulos coherentes. Si en un punto
$x\in X$ está definida la longitud geométrica $\lambda_x(\mathscr F)$ de
$\mathscr F$ en ese punto, lo mismo ocurre con las longitudes geométricas
$\lambda_x(\mathscr F')$ y $\lambda_x(\mathscr F'')$, de $\mathscr F'$ y
$\mathscr F''$ en el punto $x$, y recíprocamente~; además, se tiene
\begin{equation}
\lambda_x(\mathscr F)=\lambda_x(\mathscr F')+\lambda_x(\mathscr F'')
\tag{4.7.6.1}
\label{IV.2.4.7.6.1-fr}
\end{equation}
\end{proposition}

\begin{proof}
La proposición se deduce inmediatamente de la definición, pues
$\operatorname{long}(\mathscr F_x)$ es finita si y solo si
$\operatorname{long}(\mathscr F'_x)$ y
$\operatorname{long}(\mathscr F''_x)$ son finitas, y se tiene
$\operatorname{long}(\mathscr F_x)
=\operatorname{long}(\mathscr F'_x)+\operatorname{long}(\mathscr F''_x)$.
\end{proof}

\begin{env}[4.7.7]
\label{IV.2.4.7.7-fr}
Bajo las hipótesis de (4.7.5), los únicos puntos $x\in X$ tales que
$\mathscr F_x$ sea un $\mathcal O_x$-módulo de longitud finita \emph{no nula} son
(en virtud de (3.1.2)) \emph{los puntos maximales de}
$\operatorname{Supp}(\mathscr F)$, como se deduce de Bourbaki,
\emph{Alg. comm.}, cap.~IV, \S~2, n\textsuperscript{o}~5, cor.~2 de la
prop.~7, pues la cuestión es local sobre $X$. Se sabe que existe un
subpreesquema cerrado $Y$ que tiene $\operatorname{Supp}(\mathscr F)$ como
espacio subyacente, y un $\mathcal O_Y$-módulo coherente $\mathscr G$ tal que
$\mathscr F=j_*(\mathscr G)$, siendo $j:Y\to X$ la inyección canónica. Si $x$
es un punto maximal de $Y$, se ve que las longitudes geométricas de
$\mathscr F$ y de $\mathscr G$ en $x$ son \emph{iguales}, pues $\mathcal O_{X,x}$ y
$\mathcal O_{Y,x}$ tienen el mismo cuerpo residual. De ello se deduce la
caracterización siguiente de la longitud geométrica:
\end{env}
